\section{HD-F1-T06: the theta half-square and the index-three Weil lattice}
\label{sec:f1-integral-saturation}

Keep the principal Eisenstein component and the actual cycles of
HD-F1-S07 and T05. Put $\mathcal O=\mathbb Z[\omega]$,
$\mathfrak p=(1-\omega)$ and
$\mathscr W=W_K(A)\cap H^4(A,\mathbb Z)$. On every fibre,
\[
 \boxed{\frac{\theta^2}{2}\in\operatorname{cl}\operatorname{CH}^2(A),
 \qquad \mathfrak p\mathscr W
       \subseteq\operatorname{cl}\operatorname{CH}^2(A).}
 \tag{F1T6.1}
\]
Both assertions have actual integral signed-cycle representatives.
In particular the quotient
\[
 \mathscr W\big/
 \bigl(\mathscr W\cap\operatorname{cl}\operatorname{CH}^2(A)\bigr)
 \tag{F1T6.2}
\]
has order dividing three. This states an upper bound on the
quotient; no nonzero obstruction class is asserted.

\paragraph{The actual integral cycle maps.}
Choose an effective theta divisor representing the principal
polarization. Two general translates intersect properly, since
translation moves any prescribed component off the other divisor.
Their intersection is an effective integral cycle $D$ with
\[
             [D]=\theta^2,\qquad \deg_\theta D=24.
 \tag{F1T6.3}
\]
Here and below ``integral cycle'' permits a sum of irreducible
components with integer coefficients. It does not assert that a
particular intersection is irreducible. Let
$Z_j=[\omega^j]^*Z$ for $j=0,1,2$, where the effective cycle
$Z$ is the one constructed on this fibre in T05. Thus
\[
 [Z_j]=\frac32\theta^2+\omega^jw,
 \qquad \deg_\theta Z_j=36,
 \qquad \mathcal O w=\mathfrak p\mathscr W.
 \tag{F1T6.4}
\]
The last equality is the exact index calculation (F1T5.6).
Since $1+\omega+\omega^2=0$, the actual signed cycles
\[
 H=Z_0+Z_1+Z_2-4D,
 \qquad V=3D-Z_1-Z_2
 \tag{F1T6.5}
\]
satisfy
\[
 [H]=\frac{\theta^2}{2},\qquad [V]=w.
 \tag{F1T6.6}
\]
Indeed the first equality is
$3(3\theta^2/2)-4\theta^2=\theta^2/2$; the second is
$3\theta^2-3\theta^2-(\omega+\omega^2)w=w$.
The effective parts displayed for $H$ have degrees108 and96.
Those for $V$ have degrees72 and72. Pulling $V$ back by the
actual polarization-preserving automorphism $[\omega]$ gives
an integral representative of $\omega w$. These two classes
span $\mathcal O w=\mathfrak p\mathscr W$ over $\mathbb Z$.
This proves (F1T6.1). The lattice quotient
$\mathscr W/\mathfrak p\mathscr W$ has order3, so (F1T6.2)
is its actual quotient under the cycle-class map, proving its
stated order bound. No division in the Chow group has been used.

\paragraph{A multiplier-one extraction on the exhibited sublattice.}
For every rational Weil class $t$ with $\delta t\in\mathfrak p\mathscr W$,
write uniquely
\[
 \delta t=aw+b\omega w,
 \qquad q=a^2-ab+b^2=\frac{\delta^2T}{18},
 \qquad T=\int_A t^2.
 \tag{F1T6.7}
\]
Choose $s\in\{0,2,-2\}$ with $s\equiv a+b\pmod3$, and put
$c=(s-a-b)/3\in\mathbb Z$. The integral cycle
\[
 \Gamma=(a+c)Z_0+(b+c)Z_1+cZ_2-\frac{3s}{2}D
 \tag{F1T6.8}
\]
has class $\delta t$. Its Weil coefficient is
$(a+c)+(b+c)\omega+c\omega^2=a+b\omega$,
and its polarization coefficient is
$3(a+b+3c)/2-3s/2=0$. Every displayed coefficient is an
integer. Splitting the coefficients into positive and negative
parts gives equal degrees, each at most
\[
 d=18\bigl(|a+c|+|b+c|+|c|+|s|\bigr).
 \tag{F1T6.9}
\]
In fact this is the degree before canceling any shared components.
It follows from $\deg Z_j=36$, $\deg D=24$ and the zero
signed degree just proved.

For the degree estimate, set $c_0=-(a+b)/3$. The three
numbers $a+c_0,b+c_0,c_0$ sum to zero, so the sum of their
absolute values is twice their largest absolute value.
The identities
\[
 \begin{split}
 4q-(2a-b)^2&=3b^2,\\
 4q-(2b-a)^2&=3a^2,\\
 4q-(a+b)^2&=3(a-b)^2
 \end{split}
\]
therefore give
\[
 |a+c_0|+|b+c_0|+|c_0|\le\frac43\sqrt q.
\]
Since $c-c_0=s/3$, replacing $c_0$ by $c$ increases that
sum by at most $|s|$. Equations (F1T6.9) and $|s|\le2$
prove the explicit bound
\[
 \boxed{c_t(A)\le
   \max\{\delta,\,4\sqrt2\,\delta\sqrt T+72\}}
 \quad\text{for }\delta t\in\mathfrak p\mathscr W.
 \tag{F1T6.10}
\]
The actual multiplier is $\delta$. For classes with
$\delta t\in3\mathscr W$ the coefficient $a+b$ is divisible
by3, because reduction modulo $\mathfrak p$ sends $\omega$ to1.
Then $s=0$ and the additive72 is absent. These bounds coexist
with the uniform, all-class bound in T05; each is supported
by its own displayed integral cycle.

\paragraph{Grothendieck frameworks and source scope.}
For DEG\_L Q-DIV, (F1T6.5) and (F1T6.8) are actual maps
to signed algebraic cycle groups, formed from pullback, addition
and intersection. For the intersection map use the nonempty
open space of translation parameters on which the theta divisors
intersect properly, and shrink to a flat family before taking
its morphism to the Chow space. Addition and the unit pullbacks
then give the displayed maps on these parameter spaces.
The local relative theta line bundle is available after the
same finite level change used in F01--F04; translating its
divisor constructs the relative intersection $D$. Their
Betti and algebraic de Rham realizations give precisely
(F1T6.6), and commute with the Gauss--Manin transport already
proved for $Z$ in T05. This is a specific integral cycle
realization over the stated family in Grothendieck1966,
p.103, footnote13, with the FGA221 Hilbert/Chow maps of F02.

The identity $[\omega]^*\theta=\theta$ is the polarization
biextension identity
$\widehat{[\omega]}\lambda_\theta[\omega]=\lambda_\theta$
of F03, with SGA7I,VII2.9.5--2.9.6 and VIII3.2 as the stated
formalism. In $K_0(A)$ every cycle in (F1T6.5) is represented
by the same integral sum of the corresponding structure
sheaves. Its first nonzero character component is its
codimension-two class, and Adams $\psi^r$ multiplies that
component by $r^2$, exactly as in SGA6,V Appendix7 and F04.
These actual maps also apply to the extraction (F1T6.8).

The numerical input is the checked S07/T05 whole cycle class;
no additional algebraicity theorem is invoked. Searches of
Hexagon, arXiv and primary journal sources for the integral
Schoen/Weil lattice and divided theta-square formulations are
recorded in the source log. The elementary cycle combinations
above carry no priority claim.
